% © 2026 Ing. David Jaroš — CC BY-NC-ND 4.0
%
% This work is licensed under a Creative Commons Attribution-NonCommercial-NoDerivatives
% 4.0 International License (CC BY-NC-ND 4.0).
%
% License History: Earlier drafts (up to v0.3) were released under CC BY 4.0.
% From v0.4 onward, all material is released under CC BY-NC-ND 4.0 to protect
% the integrity of the theoretical work during ongoing academic development.
%
% See LICENSE.md for full license text.
%
% File: research_tracks/T3_ALPHA/self_consistency_fixed_point.tex
%
% Purpose: Document the self-consistent fixed-point analysis for the prime
%   attractor condition n*(B_modular(p)) = p, based on the modular volume
%   identification B(p) = (p+1)/3.
%
% Numerical results generated by: tools/self_consistency_scan.py (2026-05-10)
%
% Classification: [MC] -- motivated conjecture pending algebraic derivation of B.

% UBT-AI-PROVENANCE-BEGIN
% schema: ubt-ai-provenance/v1
% tier: C_working
% ai_assistance: disclosed
% human_review: risk-based
% editorial_responsibility: Ing. David Jaroš
% policy: ../../AI_PROVENANCE.md
% notice: Working material; exhaustive human review is not claimed.
% UBT-AI-PROVENANCE-END

\ifdefined\INCLUDEMODE
\else
  \documentclass[12pt,a4paper]{article}
  \usepackage[a4paper, margin=2.5cm]{geometry}
  \usepackage{amsmath, amssymb, amsthm}
  \usepackage{mathtools}
  \usepackage{hyperref}
  \usepackage{xcolor}
  \usepackage{booktabs}
  \usepackage{longtable}
  \usepackage{mdframed}
  \usepackage{tcolorbox}
  \tcbuselibrary{skins,breakable}

  \newtheorem{theorem}{Theorem}[section]
  \newtheorem{lemma}[theorem]{Lemma}
  \newtheorem{proposition}[theorem]{Proposition}
  \theoremstyle{definition}
  \newtheorem{definition}[theorem]{Definition}
  \theoremstyle{remark}
  \newtheorem{remark}[theorem]{Remark}

  \newmdenv[backgroundcolor=yellow!15, linecolor=orange!70, linewidth=1.5pt]{gapbox}
  \newmdenv[backgroundcolor=green!12, linecolor=green!60!black, linewidth=1.5pt]{resultbox}
  \newmdenv[backgroundcolor=blue!6, linecolor=blue!45, linewidth=1.5pt]{insightbox}
  \newmdenv[backgroundcolor=red!8, linecolor=red!60, linewidth=1.5pt]{warnbox}

  \newtcolorbox{statusbox}[1][]{colback=gray!8!white,colframe=black!60,
    arc=3pt,boxrule=1pt,fonttitle=\bfseries,title=Status Summary,#1}

  \newcommand{\Veff}{V_{\mathrm{eff}}}
  \newcommand{\Neff}{N_{\mathrm{eff}}}

  \title{\textbf{Self-Consistent Fixed-Point Equation for $\alpha^{-1}$}\\[4pt]
         \large Track T3\_ALPHA: Approach A --- $n^*(B(p)) = p$ without free parameter}
  \author{Ing.\ David Jaro\v{s}}
  \date{2026-05-10}

  \begin{document}
  \maketitle
\fi

%──────────────────────────────────────────────────────────────────────────────
\section{Motivation and Setup}
\label{sec:sc:motivation}
%──────────────────────────────────────────────────────────────────────────────

\subsection{The B-gap problem}

The UBT prime-attractor route to $\alpha^{-1}_{\mathrm{bare}} = 137$ uses the
effective potential
\begin{equation}
  \Veff(n) = n^2 - B\,n\ln n ,
  \label{eq:Veff}
\end{equation}
whose minimum satisfies the transcendental equation
\begin{equation}
  n^*(B) = \frac{B}{2}\bigl(\ln n^*(B) + 1\bigr) .
  \label{eq:nstar_transcendental}
\end{equation}
At the phenomenological value $B_{\mathrm{phenom}} \approx 46.298$, one obtains
$n^* = 137.06 \approx 137$.  However, $B_{\mathrm{phenom}}$ is \emph{not
derived} from the UBT action $S[\Theta]$.  This is Gap~G137-B.

\subsection{The modular-volume candidate for $B$}

From the arithmetic geometry of the congruence subgroup $\Gamma_0(p)$
(see \texttt{reports/gamma0\_137\_invariants.md}~\S{}4), the normalised
hyperbolic volume of the modular curve $X_0(p)$ at prime level $p$ is
\begin{equation}
  \frac{\mathrm{vol}(X_0(p))}{\pi}
  = \frac{\mu(\Gamma_0(p))}{3}
  = \frac{p+1}{3}
  \eqqcolon B_{\mathrm{mod}}(p) .
  \label{eq:Bmod}
\end{equation}
This is a purely arithmetic invariant, independent of $\alpha$.  At $p = 137$:
$B_{\mathrm{mod}}(137) = 138/3 = 46$, which reproduces $B_{\mathrm{phenom}}$
to within $0.64\%$.

\subsection{The self-consistency condition}

If UBT generates a winding-mode structure at level $n^* = p$ (prime), and if
the effective coupling $B$ equals $B_{\mathrm{mod}}(p)$, then the requirement
that $n^*(B_{\mathrm{mod}}(p)) = p$ defines a \textbf{self-consistent
fixed-point equation for $p$}:
\begin{equation}
  \boxed{p = \frac{p+1}{3} \cdot \frac{\ln p + 1}{2}}
  \label{eq:SC_basic}
\end{equation}
Rearranging:
\begin{align}
  6p &= (p+1)(\ln p + 1) \notag\\
  6p &= p\ln p + p + \ln p + 1 \notag\\
  5p - 1 &= (p+1)\ln p \notag\\
  \frac{5p-1}{p+1} &= \ln p .
  \label{eq:SC_exact}
\end{align}
This is the \textbf{exact algebraic self-consistency condition}.  It has no free
parameter.

\subsection{Refined condition including elliptic-point correction}

Including the elliptic-point correction $\nu_2(p)/4$ (where $\nu_2(p) =
1 + \left(\tfrac{-4}{p}\right)$ counts order-2 elliptic points of $\Gamma_0(p)$):
\begin{equation}
  B_{\mathrm{ref}}(p) = \frac{p+1}{3} + \frac{\nu_2(p)}{4} ,
  \label{eq:Bref}
\end{equation}
the refined self-consistency condition becomes
\begin{equation}
  p = B_{\mathrm{ref}}(p)\cdot\frac{\ln p + 1}{2} .
  \label{eq:SC_refined}
\end{equation}
At $p=137$: $\nu_2=2$, so $B_{\mathrm{ref}}(137) = 46 + 0.5 = 46.5$ (versus
$B_{\mathrm{phenom}} \approx 46.298$; error $0.44\%$).

%──────────────────────────────────────────────────────────────────────────────
\section{Numerical Scan Results}
\label{sec:sc:numerical}
%──────────────────────────────────────────────────────────────────────────────

The following results are generated by \texttt{tools/self\_consistency\_scan.py}.
The gap $\delta(p) = n^*(B_{\mathrm{mod}}(p)) - p$ measures how close each prime
$p$ is to being a self-consistent fixed point.

\subsection{Basic scan: $B_{\mathrm{mod}}(p) = (p+1)/3$}

Sorted by $|\delta(p)|$, top entries from the scan over $p \in [50, 500]$:

\begin{center}
\begin{tabular}{rrrrc}
  \toprule
  $p$ & $B_{\mathrm{mod}}$ & $n^*(B_{\mathrm{mod}})$ & $\delta(p)$ & $|\delta|<5$? \\
  \midrule
  139 & 46.667 & 138.364 & $-0.636$ & YES \\
  \textbf{137} & \textbf{46.000} & \textbf{135.989} & $-1.011$ & \textbf{YES} \\
  149 & 50.000 & 150.319 & $+1.319$ & YES \\
  151 & 50.667 & 152.726 & $+1.726$ & YES \\
  131 & 44.000 & 128.899 & $-2.101$ & YES \\
  127 & 42.667 & 124.201 & $-2.800$ & YES \\
  157 & 52.667 & 159.976 & $+2.976$ & YES \\
  163 & 54.667 & 167.269 & $+4.269$ & YES \\
  \bottomrule
\end{tabular}
\end{center}

\subsection{Refined scan: $B_{\mathrm{ref}}(p) = (p+1)/3 + \nu_2(p)/4$}

\begin{center}
\begin{tabular}{rrrrrc}
  \toprule
  $p$ & $\nu_2$ & $B_{\mathrm{ref}}$ & $n^*(B_{\mathrm{ref}})$ & $\delta_r(p)$ & $|\delta_r|<5$? \\
  \midrule
  139 & 0 & 46.667 & 138.364 & $-0.636$ & YES \\
  \textbf{137} & \textbf{2} & \textbf{46.500} & \textbf{137.770} & $+0.770$ & \textbf{YES} \\
  151 & 0 & 50.667 & 152.726 & $+1.726$ & YES \\
  131 & 0 & 44.000 & 128.899 & $-2.101$ & YES \\
  127 & 0 & 42.667 & 124.201 & $-2.800$ & YES \\
  149 & 2 & 50.500 & 152.124 & $+3.124$ & YES \\
  113 & 2 & 38.500 & 109.678 & $-3.322$ & YES \\
  \bottomrule
\end{tabular}
\end{center}

\noindent\textbf{Observation}: The refined formula brings $p=137$ to
$|\delta_r| = 0.770$ (from $1.011$ for the basic formula), passing the
$1\%$ threshold ($1\% \times 137 = 1.37$).  The closest prime is $p=139$
with $|\delta| = 0.636$ for both formulas.

\subsection{Exact algebraic condition: $(5p-1)/(p+1) = \ln p$}

\begin{center}
\begin{tabular}{rrrr}
  \toprule
  $p$ & LHS $= (5p-1)/(p+1)$ & RHS $= \ln p$ & Gap \\
  \midrule
  139 & 4.957143 & 4.934474 & $+0.022669$ \\
  \textbf{137} & \textbf{4.956522} & \textbf{4.919981} & $+0.036541$ \\
  149 & 4.960000 & 5.003946 & $-0.043946$ \\
  151 & 4.960526 & 5.017280 & $-0.056754$ \\
  \bottomrule
\end{tabular}
\end{center}

\noindent\textbf{Finding}: The exact condition $(5p-1)/(p+1) = \ln p$ does
\emph{not} have an exact prime solution.  The closest primes are $p=139$
(gap $= 0.023$) and $p=137$ (gap $= 0.037$).  The gap for $p=137$ is
$0.74\%$ of $\ln p$.

\subsection{Focus on $p = 137$}

\begin{center}
\begin{tabular}{ll}
  \toprule
  Quantity & Value \\
  \midrule
  $\nu_2(137)$ & 2 \\
  $B_{\mathrm{mod}}(137)$ & $138/3 = 46.000$ \\
  $B_{\mathrm{ref}}(137)$ & $138/3 + 2/4 = 46.500$ \\
  $B_{\mathrm{phenom}}$ (target) & $46.298$ \\
  $n^*(B_{\mathrm{mod}})$ & $135.989$; $\delta = -1.011$ ($-0.74\%$) \\
  $n^*(B_{\mathrm{ref}})$ & $137.770$; $\delta = +0.770$ ($+0.56\%$) \\
  Algebraic gap $(5p-1)/(p+1) - \ln p$ & $+0.037$ ($+0.74\%$ of $\ln p$) \\
  Passes $1\%$ threshold (basic)? & YES ($|\delta|=1.011 < 1.37$) \\
  Passes $1\%$ threshold (refined)? & YES ($|\delta|=0.770 < 1.37$) \\
  \bottomrule
\end{tabular}
\end{center}

%──────────────────────────────────────────────────────────────────────────────
\section{Analysis}
\label{sec:sc:analysis}
%──────────────────────────────────────────────────────────────────────────────

\subsection{Is $p=137$ a fixed point?}

\textbf{Approximately, yes.}  With the basic modular ansatz $B = (p+1)/3$,
the self-consistency gap is $|\delta(137)| = 1.011$, corresponding to $-0.74\%$.
With the refined formula $B_{\mathrm{ref}} = (p+1)/3 + \nu_2/4 = 46.5$, the
gap closes further to $|\delta_r(137)| = 0.770$ ($+0.56\%$).  Both are within
the $1\%$ kill condition.

However, $p=139$ has a smaller gap ($|\delta|=0.636$) for the basic formula,
and tied with $p=137$ for smallest gap at the refined level.  Thus $p=137$ is
\emph{not uniquely} selected by the self-consistency condition alone.

\subsection{What the exact algebraic condition tells us}

The exact equation $(5p-1)/(p+1) = \ln p$ arises from combining
$B_{\mathrm{mod}}(p) = (p+1)/3$ with the stationarity condition of
$\Veff$ at $n^* = p$.  Numerically, the LHS and RHS converge near
$p \approx 137$--$139$:

\begin{align}
  \text{LHS} &= \frac{5p-1}{p+1} \xrightarrow{p\to\infty} 5 , \\
  \text{RHS} &= \ln p \xrightarrow{p\to\infty} +\infty .
\end{align}
The two curves cross near $p \approx 138.5$ (non-integer, between 137 and 139).
The exact crossing is at $p_0 \approx 138.5$; neither 137 nor 139 is the exact
root, but both are the closest primes.

\begin{insightbox}
\textbf{Key result}: The self-consistency equation $(5p-1)/(p+1) = \ln p$
selects a window $p \in [137, 139]$ — a twin-prime pair — as the region where
the condition is most nearly satisfied.  The prime $p = 137$ is the lower bound
of this window, and $p = 139$ the upper.  Neither is an exact fixed point.

This twin-prime selection is structurally related to the complementary
elliptic-point structures of $\Gamma_0(137)$ and $\Gamma_0(139)$ (see
\texttt{reports/gamma0\_137\_invariants.md}~\S{}2.2).
\end{insightbox}

\subsection{Classification of the self-consistency condition}

\begin{center}
\begin{tabular}{lll}
  \toprule
  Step & Status & Reason \\
  \midrule
  $\Veff(n) = n^2 - Bn\ln n$ & [L1] & From UBT field on $S^1_\psi$ \\
  $B_{\mathrm{mod}}(p) = (p+1)/3$ & [MC] & Modular volume ansatz \\
  Self-consistency $n^*(B(p)) = p$ & [MC] & No first-principles derivation \\
  $p=137$ approximately satisfies & [MC] & Within $1\%$; not unique \\
  Exact equation has no prime root & [L0] & Numerical fact \\
  Twin-prime window $\{137,139\}$ & [MC] & Structural, needs theory \\
  \bottomrule
\end{tabular}
\end{center}

%──────────────────────────────────────────────────────────────────────────────
\section{Impact on the Alpha Track}
\label{sec:sc:impact}
%──────────────────────────────────────────────────────────────────────────────

\begin{gapbox}
\textbf{Open gap G-SC}: The self-consistency condition
$(5p-1)/(p+1) = \ln p$ \emph{approximately} selects the twin-prime window
$\{137, 139\}$ but does not uniquely select $p = 137$ over $p = 139$.
A mechanism that breaks the degeneracy — such as the elliptic-point
structure $\nu_2(137)=2, \nu_3(137)=0$ versus $\nu_2(139)=0, \nu_3(139)=2$
— is needed for a complete algebraic determination.
\end{gapbox}

The self-consistency approach has identified a structural constraint but is
insufficient alone.  It must be combined with:
\begin{enumerate}
  \item The elliptic-point correction $\nu_2/4$ which brings $B_{\mathrm{ref}}(137)$
        closer to $B_{\mathrm{phenom}}$ than $B_{\mathrm{ref}}(139)$.
  \item The Hecke eigenvalue analysis (Approach~B) which may distinguish 137 from 139.
  \item A derivation of $B_{\mathrm{mod}}(p) = (p+1)/3$ from $S[\Theta]$.
\end{enumerate}

%──────────────────────────────────────────────────────────────────────────────
\begin{statusbox}
\textbf{Výsledek}:
Nejbližší self-konzistentní prvočíslo (základní $B_{\mathrm{mod}}$):
$p = 139$ s mezerou $|\delta| = 0.636$.
$p = 137$: mezera $|\delta| = 1.011$ (základní), $|\delta_r| = 0.770$ (refinovaná).

\textbf{Je $p=137$ pevným bodem?}:
PŘIBLIŽNĚ — s přesností $0.74\%$ (základní) nebo $0.56\%$ (refinovaná).
Splňuje podmínku $|\delta| < 1\%$ pro refinovaný vzorec.
Není to \emph{přesný} pevný bod; $p=139$ je bližší základní podmínce.
Algebraická podmínka $(5p-1)/(p+1) = \ln p$ nemá přesné prvočíselné řešení;
kříží se mezi $p=137$ a $p=139$.

\textbf{Klasifikace}: [MC] — motivovaná shoda, není to algebraicky nucený výsledek.

\textbf{Dopad na alpha track}: Přístup A sám nepostačuje pro jedinečný výběr $p=137$.
Je potřeba doplňující mechanismus (Hecke struktura nebo derivace $B$ z $S[\Theta]$).
Podmínka přibližně vzata přežívá kill condition $|\delta|/p < 1\%$.
\end{statusbox}

\ifdefined\INCLUDEMODE
\else
  \end{document}
\fi
